% Rozdział 1 — Wstęp (budżet: 2,5 str.)
% Wersja dopracowana 09.11.2026 na podstawie opcji C (zatwierdzonej 08.11.2026).
% Treść jest teraz tutaj; wstep-zagajenie/wstep_opcja_C.tex nie jest już wczytywany.
% Poprzednia wersja: rozdzialy/01_wstep_backup_2026-11-09.tex.
% Zamknięte pozycje rejestru: A9 (klucze cytowań), B9/W5 (masa źródła przypisana
% źródłu), B15/W2 (brzmienie o przyrządach w misjach), B26 (cele zawężone do
% zebranych danych — do potwierdzenia przez autora).
% 09.11 wieczorem (kris): mniej cytowań (1 na tezę), U_T = 4–11 kV bez 10–18 kV.
\chapter{Wstęp}
\label{ch:wstep}

Analiza składu pierwiastkowego jest jednym z~podstawowych badań
materiałowych. W~laboratoriach wykonuje się ją najczęściej metodą
spektroskopii EDS (\angielski{Energy Dispersive X-ray Spectroscopy})
w~elektronowym mikroskopie skaningowym (SEM)~\cite{goldstein2018sem}.
Wiązka elektronów pada na~próbkę, a~zawarte w~niej pierwiastki emitują
promieniowanie X o~energiach charakterystycznych dla każdego z~nich.
Spektrometr mierzy energie fotonów, co pozwala ustalić skład próbki.

Ten sam rodzaj badań jest potrzebny w~misjach kosmicznych, w~których
określa się skład skał i~gruntu na~powierzchni Księżyca, Marsa czy
planetoid. Sprowadzanie próbek na~Ziemię jest kosztowne i~możliwe tylko
w~nielicznych misjach, dlatego analizę wykonuje się na~miejscu,
przyrządami zamontowanymi na~lądowniku lub łaziku~\cite{gellert2015apxs}.
Spektrometry APXS łazików marsjańskich wzbudzają próbkę cząstkami alfa
i~promieniowaniem X ze źródła radioizotopowego~\cite{rieder2003apxs},
a~przyrząd PIXL łazika Perseverance -- promieniowaniem z~miniaturowej
lampy rentgenowskiej~\cite{allwood2020pixl}.

Wzbudzanie próbki bezpośrednio wiązką elektronów, tak jak w~EDS, nie
zostało dotąd zastosowane w~przyrządach wykorzystanych w~misjach.
Wymaga ono działa elektronowego, które w~mikroskopie wraz z~kolumną
elektronooptyczną i~układem próżniowym ma znaczne wymiary, masę i~pobór
mocy. Projekty miniaturowych mikroskopów SEM do badań Księżyca i~Marsa
nadal opierały się na~klasycznej kolumnie elektronowej~\cite{gaskin2012mvpsem}.
Na~powierzchni ciał pozbawionych atmosfery, takich jak Księżyc czy
planetoidy, panuje natomiast naturalna próżnia, w~której wiązka
elektronów może się rozchodzić bez osobnej komory.

Zmniejszenie wymiarów źródła elektronów umożliwia technologia MEMS
(\angielski{Micro-Electro-Mechanical Systems}), w~której wytwarza się
struktury o~wymiarach od~mikrometrów do~milimetrów. Na~Politechnice
Wrocławskiej powstały w~tej technologii m.in.\ polowe źródła elektronów,
mikropompy próżniowe i~mikroskop elektronowy~\cite{krysztof2019imaging},
a~także przyrządy przeznaczone do~misji kosmicznych~\cite{dziuban2022methane}.
Opracowano tam również miniaturowe źródło promieniowania rentgenowskiego
MEMS z~zimną katodą z~nanorurek węglowych (CNT).
Zimna katoda nie wymaga elementu grzejnego, a~struktura źródła ma masę
kilku gramów~\cite{urbanski2026phd}. Źródło to zaproponowano do~budowy
kompaktowego spektrometru XRF, w~którym próbkę wzbudza promieniowanie
X~\cite{urbanski2026xrf}.

W~niniejszej pracy rozwiązanie to wykorzystano do~budowy prototypu
miniaturowego analizatora. Wiązka elektronów, ogniskowana polem
magnetycznym, pada na~target w~postaci membrany krzemowej. Wytworzone
w~niej promieniowanie X wzbudza badaną próbkę, a~spektrometr z~dyspersją
energii (EDS) rejestruje promieniowanie charakterystyczne próbki. Układ
taki stanowi miniaturowy analizator EDS w~technologii MEMS.

Tematyka pracy wpisuje się w~dyscyplinę automatyka, elektronika,
elektrotechnika i~technologie kosmiczne, ponieważ łączy elektronikę
układów zasilania źródła elektronów i~akwizycji widm, elektrotechnikę
wysokich napięć przyspieszających i~ogniskowania wiązki oraz
projektowanie przyrządu przeznaczonego do~misji kosmicznych.

Praca jest również zgodna z~profilem kierunku Inteligentna Elektronika.
Analizator EDS można traktować jako czujnik składu pierwiastkowego.
Wielkością mierzoną jest skład chemiczny próbki, a~sygnałem wyjściowym
widmo energii fotonów promieniowania X. Podobnie jak w~innych układach
czujnikowych, można w~nim wyróżnić element wzbudzający, czyli źródło
elektronów z~targetem wytwarzającym promieniowanie X, element przetwarzający, czyli detektor z~torem akwizycji,
oraz układ zasilania i~sterowania, który ustala warunki pracy źródła,
m.in.\ napięcie ekstrakcyjne i~przyspieszające.

W~docelowym przyrządzie kosmicznym taki czujnik pracowałby samodzielnie,
bez bieżącego udziału operatora, ze względu na~opóźnienia łączności
z~Ziemią. Wymaga to elektroniki, która steruje emisją elektronów,
kontroluje stabilność prądu wiązki, rejestruje widmo i~wstępnie je
przetwarza, np.~rozpoznając linie charakterystyczne pierwiastków przed
wysłaniem danych. Zagadnienia te łączą mikrosystemy MEMS z~układami
pomiarowymi i~przetwarzaniem sygnałów, co odpowiada zakresowi kierunku
Inteligentna Elektronika.

\section{Cel i~zakres pracy}
\label{sec:cel}

Celem pracy jest opracowanie, wykonanie oraz badania prototypu
miniaturowego analizatora EDS w~technologii MEMS, przeznaczonego
do~zastosowań kosmicznych. Zakres pracy obejmuje część inżynierską
i~badawczą.

Część inżynierska:
\begin{itemize}
  \item zaprojektowanie elementów źródła elektronów (emitera, bramki
    i~targetu) oraz masek do~procesu technologicznego,
  \item wykonanie struktur krzemowych i~szklanych, naniesienie kompozytu
    CNT-PVP na~emiter i~połączenie elementów bondingiem anodowym,
  \item zestawienie stanowiska pomiarowego z~komorą próżniową, zasilaniem
    wysokim napięciem i~spektrometrem rentgenowskim.
\end{itemize}

Część badawcza:
\begin{itemize}
  \item pomiary emisji polowej emitera CNT, czyli zależności prądu emisji
    od~napięcia,
  \item ocena wpływu magnesu ogniskującego na~kształt wiązki obserwowanej
    na~ekranie luminescencyjnym,
  \item rejestracja widma promieniowania X wzbudzanego wiązką elektronów
    i~identyfikacja linii charakterystycznych,
  \item ocena możliwości zastosowania technologii MEMS do~miniaturyzacji
    analizatorów EDS przeznaczonych do~zastosowań kosmicznych.
\end{itemize}

% B26: z celów usunięto wpływ geometrii bramki, średnicę wiązki w µm
% i porównanie liczby fotonów z magnesem i bez, bo brak takich danych.
% Jeśli autor je dostarczy, przywrócić odpowiednie punkty.
Badania prowadzono w~komorze próżniowej przy ciśnieniu rzędu
\SI{1e-3}{\pascal} (\SI{1e-5}{\hecto\pascal}). Elementy połączono
bondingiem anodowym w~atmosferze, dlatego struktura nie jest hermetyczna.
Hermetyzacja z~mikropompą, istotna w~docelowym przyrządzie, nie wchodzi
w~zakres pracy, podobnie jak konstrukcja zasilacza wysokiego napięcia
i~detektora. Widma rejestrowano przy napięciu przyspieszającym $U_\mathrm{T}$
w~zakresie \SIrange{4}{11}{\kilo\volt}.
